\documentclass[11pt,a4paper]{article}
\usepackage[T1]{fontenc}
\usepackage[utf8]{inputenc}
\usepackage[margin=24mm]{geometry}
\usepackage{amsmath,amssymb,array,booktabs,longtable}
\usepackage[colorlinks=true,linkcolor=blue,urlcolor=blue]{hyperref}
\newcommand{\mat}[1]{\mathbf{#1}}
\newcommand{\code}[1]{\texttt{\detokenize{#1}}}
\newcommand{\src}[1]{\href{../../#1}{\nolinkurl{#1}}}
\newcommand{\prereq}[1]{\textit{Builds on: #1}\par}
\newcommand{\checkpoint}[1]{\par\textbf{Check your understanding.} #1\par}
\newcommand{\topic}[2]{\boxed{\text{\hyperref[#1]{#2}}}}
\DeclareMathOperator{\tr}{tr}
\DeclareMathOperator{\diag}{diag}
\DeclareMathOperator{\cof}{cof}
\DeclareMathOperator{\Div}{Div}
\setlength{\parindent}{0pt}
\setlength{\parskip}{6pt}
\setlength{\emergencystretch}{2em}
\renewcommand{\arraystretch}{1.18}
\numberwithin{equation}{section}
\title{Mathematical Foundations and Derivations\\for CK-MPM and Sand Simulation}
\author{A learning companion to the paper-to-code map}
\date{}
\begin{document}
\maketitle

\textbf{Purpose.} Start with derivatives and small matrices, then follow how
geometry, continuum mechanics and discretization produce the equations used
in CK-MPM. This is a connected learning route for the paper's core numerical
method and the current sand model, not a claim to prove every theorem or GPU
optimization in the paper. The companion
\href{paper_code_map.pdf}{\emph{Paper-to-Code Guide}} locates implementation
functions. Equations numbered here belong to this document; ``paper Eq.''
refers to Liu, Wang and Li's CK-MPM paper~\cite{ck}.

The examples and derivations below use ordinary column-vector notation.
A result derived from stated assumptions, a model choice and an expression
copied from the author implementation are different kinds of evidence.
Sections~\ref{sec:mls} and~\ref{sec:history} identify places where that distinction
matters. No simulation code is changed by this guide.

\tableofcontents
\newpage

\section{Dependency map and notation}\label{sec:map}
The boxes are clickable reading routes. The table gives the cross-branch
prerequisites that a simple chain cannot show.
\[
\topic{sec:jacobian}{Derivatives/Jacobian}\to
\topic{sec:deformation}{Deformation}\to
\topic{sec:volume}{Volume and mass}
\]
\[
\topic{sec:nanson}{Area/Nanson}\to
\topic{sec:stress}{Stress conversion}\to
\topic{sec:weak}{Weak form}\to\topic{sec:mpm}{MPM}
\]
\[
\topic{sec:kernel}{Kernel}\to\topic{sec:dual}{Dual-grid moments}\to
\topic{sec:apic}{Affine fit}\to\topic{sec:mls}{MLS forces}
\]
\[
\topic{sec:svd}{SVD}\to\topic{sec:energy}{Elastic stress}\to
\topic{sec:plastic}{Yield/return map}\to\topic{sec:trace}{One timestep}
\]
\begin{longtable}{>{\raggedright\arraybackslash}p{.29\linewidth}>{\raggedright\arraybackslash}p{.62\linewidth}}
\toprule Result & Immediate prerequisites \\\midrule\endhead
Deformation F & Derivative, Jacobian, first-order Taylor approximation. \\
Volume J and density & F, determinant, change of variables, conservation of mass. \\
Nanson's formula & F, oriented area/cross product, determinant and inverse transpose. \\
Piola stress P & Cauchy traction, equality of physical force, Nanson's formula. \\
Weak internal force & Momentum balance, integration by parts, P, chain rule. \\
Particle/grid MPM force & Weak form, shape functions, particle quadrature and lumped mass. \\
Deformation evolution & Chain rule in space and time; velocity gradient. \\
CK first-order consistency & Compact weights, partition of unity, two grid offsets. \\
APIC affine fit & Outer products, weighted least squares, dual-grid moment identities. \\
MLS force approximation & Affine fit, shape gradients, the weak internal force. \\
Hencky stress & SVD, logarithmic stretch, differentiating elastic energy. \\
Sand return mapping & Hencky stress, volumetric/deviatoric split and a yield criterion. \\\bottomrule
\end{longtable}

\textbf{Notation.} $\mat{X}$ is a material/reference position and $\mat{x}$ a
current/spatial position. Roman $i,j$ index spatial components, capital $A,B$
index reference components; repeated component indices are summed. A particle
index $p$ and a grid index $k$ are not summed unless a summation is written.
$\nabla f$ is a column gradient, $(\nabla\mat{v})_{ij}=\partial v_i/\partial x_j$,
and $(\Div\mat{P})_i=\partial P_{iA}/\partial X_A$.

Until plasticity is introduced, F is the Jacobian of the full deformation map.
In the plasticity sections we distinguish total F from elastic $\mat{F}^e$.
A corrected elastic matrix stored by a solver must not automatically be treated
as the full geometric deformation Jacobian.

\section{Derivatives, matrices and the Jacobian}\label{sec:jacobian}
\prereq{Functions, coordinates, addition and multiplication.}
A derivative describes how an output changes under a small input change:
\[
f'(a)=\lim_{h\to0}\frac{f(a+h)-f(a)}h,
\qquad f(a+h)=f(a)+f'(a)h+O(h^2)
\]
for a sufficiently smooth function. For $f(x_1,x_2)$, the partial derivative
$\partial f/\partial x_1$ varies $x_1$ while holding $x_2$ fixed. Collecting the
partial derivatives gives the gradient, so a small vector increment obeys
$f(\mat{x}+\mat{h})\approx f(\mat{x})+\nabla f(\mat{x})^T\mat{h}$.

For a vector-valued function $\mat{f}:\mathbb R^n\to\mathbb R^m$, each output
has its own gradient. Stacking those gradients as rows gives its
\emph{Jacobian matrix}:
\begin{equation}\label{eq:jacobian}
\mat{J}_{\mat f}=\left[\frac{\partial f_i}{\partial x_j}\right],
\qquad
\mat f(\mat x+\mat h)\approx\mat f(\mat x)+\mat J_{\mat f}(\mat x)\mat h.
\end{equation}
For $\mat f(X_1,X_2)=(2X_1+X_2,X_2^2)$, the Jacobian at $(1,1)$ is
$\left[\begin{smallmatrix}2&1\\0&2\end{smallmatrix}\right]$.
Thus $(.01,0)$ changes the output by approximately $(.02,0)$, whereas
$(0,.01)$ changes it by approximately $(.01,.02)$.

A transpose exchanges rows and columns. For column vectors $\mat a,\mat b$,
$\mat a^T\mat b$ is a scalar dot product, while $\mat a\mat b^T$ is an outer
product with entry $a_i b_j$. For matrices, define the double contraction
$\mat A:\mat B=\sum_{ij}A_{ij}B_{ij}=\tr(\mat A^T\mat B)$.
Matrix order matters: generally $\mat A\mat B\ne\mat B\mat A$.

The chain rule is matrix multiplication:
$\mat J_{\mat g\circ\mat f}=\mat J_{\mat g}\mat J_{\mat f}$.
An inverse matrix reverses an invertible local map. It does not mean taking
the reciprocal of every matrix entry.
\checkpoint{The Jacobian is a matrix; its determinant is a scalar. The symbol J
later denotes the determinant of F, not the Jacobian matrix itself.}

\section{Deformation and the velocity gradient}\label{sec:deformation}
\prereq{Section~\ref{sec:jacobian}.}
A motion map sends each material point to its current position:
\[
\mat x=\boldsymbol\phi(\mat X,t),\qquad
F_{iA}=\frac{\partial\phi_i}{\partial X_A}.
\]
Applying the Jacobian approximation to neighboring material points gives
\begin{equation}\label{eq:line-map}
\mathrm d\mat x=\mat F\,\mathrm d\mat X.
\end{equation}
F records local stretching, shearing and rotation. Translation disappears
under differentiation. For displacement $\mat u=\boldsymbol\phi-\mat X$,
$\mat F=\mat I+\nabla_X\mat u$.

Examples: a translation has $\mat F=\mat I$; a rigid rotation has
$\mat F=\mat R$, $\mat R^T\mat R=\mat I$; and a uniform stretch has
$\mat F=\diag(s_1,s_2,s_3)$. The shear map $x_1=X_1+\gamma X_2$, $x_2=X_2$
has $\mat F=\left[\begin{smallmatrix}1&\gamma\\0&1\end{smallmatrix}\right]$.

Let $\mat v(\mat x,t)=\partial_t\boldsymbol\phi(\mat X,t)$ evaluated along
that motion. Differentiating F and using the chain rule gives
\begin{equation}\label{eq:Fdot}
\dot F_{iA}=\frac{\partial v_i}{\partial x_j}F_{jA},
\qquad \dot{\mat F}=\mat L\mat F,\quad\mat L=\nabla_x\mat v.
\end{equation}
Forward Euler over a small step gives
$\mat F^{\rm trial}=(\mat I+\Delta t\,\mat L)\mat F^n$.
This is the deformation-update structure in paper Eq.~(19). Later the affine
fit approximates L.

For a scalar spatial function $w(\mat x)$, differentiation of
$w(\boldsymbol\phi(\mat X))$ gives another useful chain rule:
\begin{equation}\label{eq:gradient-transform}
\nabla_X w=\mat F^T\nabla_x w.
\end{equation}
\checkpoint{Velocity has units length/time, L has units $1$/time, and F is
dimensionless. Therefore $\mat I+\Delta t\,\mat L$ has consistent units.}

\section{Determinants, volume and mass}\label{sec:volume}
\prereq{Section~\ref{sec:deformation}; dot and cross products.}
Three short edge vectors $\mat a,\mat b,\mat c$ span an oriented volume
$\mat a\cdot(\mat b\times\mat c)$. A matrix transforms that triple product by
its determinant:
\[
(\mat F\mat a)\cdot[(\mat F\mat b)\times(\mat F\mat c)]
=\det(\mat F)\,\mat a\cdot(\mat b\times\mat c).
\]
Writing $J=\det\mat F$, an orientation-preserving deformation has $J>0$ and
\begin{equation}\label{eq:volume}
\mathrm dv=J\,\mathrm dV.
\end{equation}
Lowercase $\mathrm dv$ is current volume; uppercase $\mathrm dV$ is reference
volume. With $\mat F=\diag(2,1,.5)$, J is one: volume is unchanged despite
changes in shape. A rotation has J one too.

A material element keeps its mass. If $\rho_R$ is reference mass density and
$\rho$ is current density, then
\begin{equation}\label{eq:mass}
\rho\,\mathrm dv=\rho_R\,\mathrm dV
\quad\Longrightarrow\quad \rho J=\rho_R.
\end{equation}
This is the reference/current mass-balance connection behind the paper's
continuum preliminaries. Reference quadrature volume $V_{p,0}$ stays fixed;
current geometric volume would be $J_pV_{p,0}$.

For an invertible F, differentiation of a determinant yields
$\dot J=J\tr(\mat F^{-1}\dot{\mat F})$.
Using \eqref{eq:Fdot} and cyclic invariance of trace gives
$\dot J=J\tr\mat L=J\nabla\cdot\mat v$.
Differentiating \eqref{eq:mass} then gives
$\dot\rho+\rho\nabla\cdot\mat v=0$.
\checkpoint{J describes local volume change, not a particle count and not the
simulation timestep. These meanings must remain separate.}

\section{Oriented area and Nanson's formula}\label{sec:nanson}
\prereq{Sections~\ref{sec:deformation} and~\ref{sec:volume}.}
A surface patch is spanned by two small tangent vectors $\mat a,\mat b$.
Their cross product contains both orientation and area:
$\mat a\times\mat b=\mat N\,\mathrm dA$, where N is the reference unit normal.
After deformation, the tangents become F a and F b, so
$(\mat F\mat a)\times(\mat F\mat b)=\mat n\,\mathrm da$.

To find the transformation, test the deformed area vector against an arbitrary
third vector $\mat F\mat c$. The determinant identity gives
\begin{align*}
[(\mat F\mat a)\times(\mat F\mat b)]\cdot\mat F\mat c
&=J(\mat a\times\mat b)\cdot\mat c,\\
\mat F^T[(\mat F\mat a)\times(\mat F\mat b)]
&=J(\mat a\times\mat b).
\end{align*}
The second equality follows because the first holds for every c. Multiplying
by $\mat F^{-T}$ proves the area-vector transformation:
\begin{equation}\label{eq:nanson}
\boxed{\mat n\,\mathrm da=J\mat F^{-T}\mat N\,\mathrm dA
=\cof(\mat F)\mat N\,\mathrm dA.}
\end{equation}
Here $\mat F^{-T}=(\mat F^{-1})^T$ and
$\cof(\mat F)=J\mat F^{-T}$ for invertible F. This is Nanson's formula;
continuum-mechanics references give this bridge explicitly~\cite{sandia,nptel}.

Taking magnitude and direction separately gives
\[
\frac{\mathrm da}{\mathrm dA}=J\|\mat F^{-T}\mat N\|,
\qquad
\mat n=\frac{\mat F^{-T}\mat N}{\|\mat F^{-T}\mat N\|}.
\]
\textbf{Worked check.} With $\mat F=\diag(2,1,.5)$ and $\mat N=\mat e_1$,
J is one but the patch area scales by $1/2$. Its two tangents stretch by one
and $1/2$. A normal along $\mat e_3$ gives area factor two instead.
\checkpoint{Why not multiply surface area by J alone? Different orientations
stretch differently. J transforms volume; Nanson transforms oriented area.}

\section{Traction, Cauchy stress and Piola stress}\label{sec:stress}
\prereq{Section~\ref{sec:nanson}.}
Traction is force per unit current area on a plane with unit normal n.
The Cauchy stress tensor $\boldsymbol\sigma$ maps that normal to traction:
$\mat t=\boldsymbol\sigma\mat n$.
The physical differential force must be identical when described per reference
area. Define first Piola stress P by $\mathrm d\mat f=\mat P\mat N\,\mathrm dA$.
Insert Nanson's formula into the current-area description:
\[
\mathrm d\mat f=\boldsymbol\sigma\mat n\,\mathrm da
=J\boldsymbol\sigma\mat F^{-T}\mat N\,\mathrm dA.
\]
Since this holds for every reference normal,
\begin{equation}\label{eq:piola}
\boxed{\mat P=J\boldsymbol\sigma\mat F^{-T}},\qquad
\boxed{\mat T=\mat P\mat F^T=J\boldsymbol\sigma}.
\end{equation}
T is the Kirchhoff stress in this notation. Sandia calls its Cauchy stress T;
we use $\boldsymbol\sigma$ to avoid confusing it with the stored
$\mat P\mat F^T$ in our shader. A symbol's name is not its physical definition.
See also the stress-measure teaching reference~\cite{alberta}.

Under ordinary balance without couple stresses, $\boldsymbol\sigma$ and T are
symmetric. P need not be symmetric because its input normal is reference-based
and its output force is spatial. It is a two-configuration tensor.

\textbf{Worked check.} If $\mat F=2\mat I$, then $J=8$ and
$\mat P=4\boldsymbol\sigma$. A reference patch grows to four times its area,
so this correctly preserves its physical force.
\checkpoint{Nanson's formula explains the conversion of stress measures. It
is not an extra sand force or an algorithm called once per particle.}

\section{Momentum balance and the weak internal force}\label{sec:weak}
\prereq{Sections~\ref{sec:stress} and~\ref{sec:deformation}; integration.}
For a material region, Newton's law equates momentum change to body and surface
forces. Applying the divergence theorem to the surface traction gives the
reference strong form
\begin{equation}\label{eq:strong}
\rho_R\ddot{\boldsymbol\phi}=\Div_X\mat P+\rho_R\mat g.
\end{equation}
The divergence theorem here means integrating a derivative inside a region is
equivalent to integrating its normal flux on the boundary. For our row-divergence
convention, $\int_{\Omega_0}\Div_X\mat P\,\mathrm dV
=\int_{\partial\Omega_0}\mat P\mat N\,\mathrm dA$.

Multiply \eqref{eq:strong} by a scalar test/shape function $w_i(\mat x)$ and
integrate. Component-wise integration by parts gives
\begin{align}
\int_{\Omega_0}\rho_R w_i\ddot\phi_a\,\mathrm dV
&=-\int_{\Omega_0}P_{aA}\frac{\partial w_i}{\partial X_A}\,\mathrm dV\notag\\
&\quad+\int_{\partial\Omega_0}w_i(PN)_a\,\mathrm dA
+\int_{\Omega_0}\rho_R w_i g_a\,\mathrm dV.\label{eq:weak}
\end{align}
Here i labels a grid shape function, while a labels a vector component: they
are independent indices. The minus sign is the integration-by-parts sign.
For an interior node or zero prescribed traction, the explicit boundary term
vanishes. Boundary conditions otherwise contribute force or constraints.

Use \eqref{eq:gradient-transform} to move the shape gradient into spatial
coordinates:
\begin{equation}\label{eq:internal-integral}
\mat f_i^{\rm int}
=-\int_{\Omega_0}\mat P\mat F^T\nabla_x w_i\,\mathrm dV
=-\int_{\Omega_0}\mat T\nabla_x w_i\,\mathrm dV.
\end{equation}
This explains both the negative sign and the $\mat P\mat F^T$ in paper Eq.~(2).
The factor F transpose is supplied by the chain rule; Nanson supplied the
stress conversion earlier in the dependency chain.

There is a second check using current volume:
$\boldsymbol\sigma\,\mathrm dv
=\boldsymbol\sigma J\,\mathrm dV=\mat T\,\mathrm dV$.
Thus $-\int_{\Omega_t}\boldsymbol\sigma\nabla w_i\,\mathrm dv$ is the same force.
\checkpoint{Do not multiply by J a second time when already using reference
volume and $\mat T=\mat P\mat F^T$. That would count the volume conversion twice.}

\section{From integrals to MPM particles and grid nodes}\label{sec:mpm}
\prereq{Section~\ref{sec:weak}; weighted sums.}
Particle quadrature approximates a reference-volume integral by
$\int_{\Omega_0}q\,\mathrm dV\approx\sum_p V_{p,0}q_p$.
Particles carry mass $m_p=\rho_{R,p}V_{p,0}$ and state history; grid nodes provide
locations for collecting and updating motion.

Interpolate a grid field by shape functions $w_{i,p}=N_i(\mat x_p)$.
The consistent mass matrix is approximated by
$M_{ij}\approx\sum_p m_p w_{i,p}w_{j,p}$.
With partition of unity $\sum_jw_{j,p}=1$, row-summing gives a diagonal
\emph{lumped} nodal mass $m_i=\sum_p m_pw_{i,p}$. This replaces a coupled mass
solve with independent nodal updates.

For a single grid, the elementary transfers and force are
\begin{align}
m_i&=\sum_p w_{i,p}m_p,&
\mat q_i&=\sum_p w_{i,p}m_p\mat v_p,\label{eq:p2g-basic}\\
\mat f_i^{\rm int}&=-\sum_p V_{p,0}\mat T_p\nabla w_{i,p},&
\mat u_i&=\frac{\mat q_i+\Delta t\,\mat f_i^{\rm int}}{m_i}
                 +\Delta t\,\mat g.\label{eq:grid-basic}
\end{align}
The divisions apply only to sufficiently massive active nodes. Particles then
gather grid velocity and advect using the new velocity. In a traditional
kernel-gradient formulation, $\mat L_p\approx\sum_i\mat u_i(\nabla w_{i,p})^T$
feeds \eqref{eq:Fdot}. The affine formulation below supplies another estimate.

In our code, P2G adds stress momentum directly to q, so the grid update only
normalizes the already-forced momentum, adds gravity and handles boundaries.
This rearranges stages; it does not justify dropping the stress force.
For CK, each grid gets its own full mass/momentum scatter, and the gathering
and physical conservation totals account for averaging the two grids.

\textbf{Tiny transfer example.} Two equal-mass particles send weight $1/2$
to the same node with velocities $+1$ and $-1$. That node has positive mass
and zero momentum, hence zero velocity. Repeating weighted averaging can erase
velocity variation. Particle count is not the same thing as independent grain
contacts: this is a discretized continuum.
\checkpoint{Forces have units mass times length/time squared. In
$V_{p,0}\mat T\nabla w$, volume contributes length cubed, stress force/area and
gradient inverse length, giving force.}

\section{Time integration, stability and the two clocks}\label{sec:time}
\prereq{Sections~\ref{sec:deformation} and~\ref{sec:mpm}.}
The particle update $\mat x^{n+1}=\mat x^n+\Delta t\,\mat v^{n+1}$ uses the
updated velocity. Together with an explicit velocity kick this is the
symplectic-Euler pattern. It does not allow an arbitrary timestep.

For a spring $\ddot x=-\omega^2x$, the same pattern gives
\[
\begin{pmatrix}x^{n+1}\\v^{n+1}\end{pmatrix}
=\begin{pmatrix}1-h^2\omega^2&h\\-h\omega^2&1\end{pmatrix}
\begin{pmatrix}x^n\\v^n\end{pmatrix}.
\]
The matrix has determinant one and trace $2-h^2\omega^2$. Oscillatory bounded
modes require $0<h\omega<2$; at the endpoint the repeated eigenvalue can be
defective. Larger steps can amplify errors even when the intended motion is
slow. Elastic waves give a related scale $h\sim\Delta x/c$.

For small isotropic elastic disturbances, define
\[
\mu=\frac{E}{2(1+\nu)},\qquad
\lambda=\frac{E\nu}{(1+\nu)(1-2\nu)},\qquad
c=\sqrt{\frac{\lambda+2\mu}{\rho}}.
\]
Our automatic acoustic estimate is $\Delta t={\rm CFL}\,\Delta x/c$.
It is a practical estimate, not a proof of nonlinear elastoplastic stability.

A numerical step advances simulated time by $\Delta t$ regardless of how long
the computer takes. With s steps per frame and R rendered frames per real
second, the observed time ratio is approximately $s\Delta t R$.
For $s=3$, $\Delta t=.0004$ and $R=60$, the ratio is .072.
Larger s adds work and may reduce R; larger $\Delta t$ changes integration
accuracy and stability. These are different controls.
\checkpoint{A 16 ms computer work budget is not automatically 16 ms of stable
physical evolution. Deriving $\Delta t=\text{desired simulated frame time}/s$
is possible, but its numerical stability must be checked separately.}

\section{Shape functions and the compact kernel}\label{sec:kernel}
\prereq{Sections~\ref{sec:jacobian} and~\ref{sec:mpm}.}
A shape function is a position-dependent interpolation weight. Useful properties
include locality, nonnegativity and partition of unity. It is not simply the
Euclidean distance; distance is an input to its formula.

The 1D linear tent is $K_{\rm lin}(r)=\max(1-|r|,0)$. Its slope jumps at zero
and at the support edges. A force that uses its derivative can therefore change
abruptly when a particle crosses a cell boundary.

The CK paper's Eq.~(4) chooses
\begin{equation}\label{eq:ck}
K(r)=\begin{cases}
1-|r|+\dfrac{\sin(2\pi|r|)}{2\pi},& |r|\le1,\\
0,&|r|>1.
\end{cases}
\end{equation}
One can motivate the coefficient without assuming the result. Start on
$0<r<1$ with $1-r+a\sin(2\pi r)$. Requiring the derivative to vanish at zero
gives $-1+2\pi a=0$, hence $a=1/(2\pi)$. This is the simple chosen correction;
it is not a proof that no other smooth compact kernels exist.

For $0<r<1$, differentiation gives
\[
K'(r)=-1+\cos(2\pi r),\qquad K''(r)=-2\pi\sin(2\pi r).
\]
The first and second derivatives approach zero at the join points. Together
with even symmetry and the zero exterior, this makes K $C^2$ continuous.
Since $K'\le0$ on $(0,1)$ and $K(0)=1$, $K(1)=0$, weights are nonnegative.
The sine integrates to zero over $(0,1)$, so symmetry gives
$\int_{\mathbb R}K(r)\,\mathrm dr=1$.

For $d\in[0,1]$,
\begin{equation}\label{eq:unity}
K(d)+K(1-d)=1,
\end{equation}
because $\sin(2\pi(1-d))=-\sin(2\pi d)$.
This gives the two cell weights $w_0=K(d)$, $w_1=1-w_0$ in
\src{shaders/interpolation.glsl}. In 3D, multiply the three axis weights;
each grid contributes $2^3=8$ nodes. The paper's Eq.~(5) is this tensor product.

\checkpoint{At $d=.25$, $w_0=.75+1/(2\pi)\approx.9092$ and $w_1\approx.0908$.
They sum to one but are not the linear weights .75 and .25.}

\section{Why two shifted grids are needed}\label{sec:dual}
\prereq{Section~\ref{sec:kernel}; weighted averages and outer products.}
Partition of unity reproduces constants. Reproducing a linear position field
also needs the first moment to agree with particle position.
On one grid, write $x_p/\Delta x=b+d$, where b is the left node coordinate and
$0\le d<1$. The weighted node coordinate is
\[
bw_0+(b+1)w_1=b+d-\frac{\sin(2\pi d)}{2\pi}.
\]
The last term is an error in linear reproduction, even though the weights sum
to one.

Use grids shifted by $-\Delta x/4$ and $+\Delta x/4$. Their local d values
differ by half a cell modulo one. Since
$\sin(2\pi(d+1/2))=-\sin(2\pi d)$, averaging the two first moments cancels
the sine errors. Tensor-product weights extend this argument to each component
in 3D. This is the key consistency role of the dual-grid construction in
paper \S4.2.

Define the combined weighted average
\[
\langle a\rangle_p=\frac12\sum_{k\in\{-1,+1\}}\sum_i w_{i,p,k}a_{i,k},
\qquad \mat r_{i,p,k}=\mat x_{i,k}-\mat x_p.
\]
For a complete interior support,
\begin{equation}\label{eq:moments}
\langle1\rangle_p=1,\qquad \langle\mat r\rangle_p=\mat0,
\qquad \mat D_p=\langle\mat r\mat r^T\rangle_p.
\end{equation}
For any vector z,
$\mat z^T\mat D\mat z=\langle(\mat z\cdot\mat r)^2\rangle\ge0$.
Thus D is positive semidefinite, and it is invertible if the positively
weighted offsets span all spatial directions. It has units length squared.

Each grid receives the full scattered mass. Summing both grid masses would
count material twice; a physical total uses half their sum. Likewise the
velocity and moment reconstruction use the combined average above.

Boundary clipping can invalidate the complete-support identities. Our GPU
checks therefore require transfer stencils to stay inside the grid. Applying
a wall constraint to velocity also exerts a physical boundary impulse; momentum
need not remain constant when gravity or walls act.
\checkpoint{Sixteen grid nodes are not sixteen new sand grains. They are two
sets of grid samples supporting the same material particle.}

\section{Deriving the affine least-squares fit}\label{sec:apic}
\prereq{Sections~\ref{sec:jacobian} and~\ref{sec:dual}.}
Approximate a nearby velocity field by
$\mat v(\mat x_p+\mat r)\approx\mat v_p+\mat C_p\mat r$.
C describes spatial variation: a skew-symmetric C is rotation, a symmetric
part describes stretching/shearing. APIC carries this variation through
particle/grid transfer~\cite{apic}.

Choose v and C to minimize the combined weighted residual
\begin{equation}\label{eq:fit}
E(\mat v,\mat C)=\tfrac12\langle\|\mat u_i-\mat v-\mat C\mat r_i\|^2\rangle.
\end{equation}
Differentiating with respect to v gives
$\mat v=\langle\mat u\rangle-\mat C\langle\mat r\rangle
=\langle\mat u\rangle$.
Differentiating entry $C_{ab}$ gives
\[
0=-\langle(u_a-v_a-C_{aj}r_j)r_b\rangle.
\]
Using the zero first moment leaves the normal equations
\begin{equation}\label{eq:affine-fit}
\mat C\mat D=\mat B,\qquad
\mat B=\langle\mat u\mat r^T\rangle,\qquad
\boxed{\mat C=\mat B\mat D^{-1}}.
\end{equation}
These are the dual-grid D and B behind paper \S4.4 and Eqs.~(20)--(21).
B has units length squared/time, so C has units inverse time.

G2P uses offsets $\widehat{\mat r}=\mat r/\Delta x$ in cell units. With
$\widehat{\mat B}=\sum w\mat u\widehat{\mat r}^{T}$ and
$\widehat{\mat D}=\sum w\widehat{\mat r}\widehat{\mat r}^{T}$, it therefore computes
\[
\mat C=\frac1{\Delta x}\widehat{\mat B}\widehat{\mat D}^{-1}.
\]
The omitted halves cancel. Forgetting the division by $\Delta x$ would change
the physical units and the amount of deformation.

P2G then transfers $w m_p(\mat v_p+\mat C_p\mat r_i)$.
For an exact affine field $\mat u_i=\mat a+\mat A\mat x_i$, the moment
identities give $\mat v_p=\mat a+\mat A\mat x_p$ and $\mat C_p=\mat A$.
This supplies an analytic test independent of another implementation.
For rigid rotation around z,
$\mat A=\left[\begin{smallmatrix}0&-\omega&0\\\omega&0&0\\0&0&0\end{smallmatrix}\right]$.
The fitted field correctly sends opposite sides in opposite directions.

\subsection{Where angular information is stored}
The grid's angular momentum contribution of one affine particle expands to
\[
\langle(\mat x_p+\mat r)\times m_p(\mat v_p+\mat C_p\mat r)\rangle
=\mat x_p\times m_p\mat v_p
 +m_p\langle\mat r\times\mat C_p\mat r\rangle.
\]
The cross terms vanish by $\langle\mat r\rangle=0$. The second term is local
spin information absent from a point-position/point-velocity pair. It depends
on $\mat C\mat D=\mat B$. Measuring APIC angular momentum using only
$\sum_p\mat x_p\times m_p\mat v_p$ omits this part.
This expansion explains the mechanism, not a full floating-point conservation
proof for our force/boundary/plasticity pipeline.

\checkpoint{For a constant velocity field, B is zero because the first moment
is zero, hence C is zero. For an affine field, the fit should recover C exactly
in exact arithmetic with a complete support.}

\section{MLS reconstruction and a convention audit}\label{sec:mls}
\prereq{Sections~\ref{sec:weak} and~\ref{sec:apic}.}
For a scalar field, fit $a_i\approx a_0+\mat q^T\mat r_i$ using the same
weighted objective. The normal equations give
$a_0=\langle a\rangle$ and
$\mat q=\mat D^{-1}\langle\mat r a\rangle$.
Reconstruct at a nearby point z while holding the fit center $\mat x_p$ fixed:
\begin{align}
a(\mat z)&\approx\tfrac12\sum_{k,i}N_{i,k}(\mat z)a_{i,k},\\
N_{i,k}(\mat z)&=w_{i,p,k}
 [1+(\mat z-\mat x_p)^T\mat D_p^{-1}\mat r_{i,p,k}],\label{eq:mls-shape}\\
\nabla_zN_{i,k}(\mat x_p)&=w_{i,p,k}\mat D_p^{-1}\mat r_{i,p,k}.\label{eq:mls-gradient}
\end{align}
The weights and moment matrix are frozen while differentiating in z. Moving
the fit center and differentiating its weights is a different operation.
The block moment matrix for basis $(1,\mat r)$ is
$\left[\begin{smallmatrix}1&0\\0&\mat D\end{smallmatrix}\right]$ because of
\eqref{eq:moments}. This is the concrete linear-basis connection to the general
MLS discussion in paper \S4.5~\cite{mls}.

Substituting \eqref{eq:mls-gradient} into \eqref{eq:internal-integral} and using
particle quadrature gives, in the column-vector convention of this guide,
\begin{equation}\label{eq:mls-force}
\boxed{\mat f_{i,k}^{\rm int}
=-\sum_p V_{p,0}w_{i,p,k}\mat T_p\mat D_p^{-1}\mat r_{i,p,k}.}
\end{equation}
For dual-grid nodal variables, the half in the physical averaged reconstruction
appears in both nodal mass and force. It cancels in the update when using the
code's full per-grid mass storage. It does not change matrix multiplication
order.

\subsection{Keep the derived expression and source expression distinct}
\label{sec:order-audit}
The current shader and the inspected author CUDA source instead form
\[
\mat A_{\rm source}=m_p\mat C_p
 -\Delta t\,V_{p,0}\mat D_p^{-1}\mat T_p,
\qquad
\mat q_{i,k}\mathrel{+}=w(m_p\mat v_p+\mat A_{\rm source}\mat r).
\]
The source puts $\mat D^{-1}$ on the left of T, while
\eqref{eq:mls-force} places it on the right. These matrices are generally
not interchangeable. For instance, in a 2D block,
\[
\mat D=\begin{pmatrix}2&0\\0&1\end{pmatrix},\quad
\mat T=\begin{pmatrix}1&1\\1&2\end{pmatrix},\quad
\mat D^{-1}\mat T=\begin{pmatrix}.5&.5\\1&2\end{pmatrix}
\ne
\mat T\mat D^{-1}=\begin{pmatrix}.5&1\\.5&2\end{pmatrix}.
\]
Symmetry of both factors does not make their product commute. The distinction
can disappear for a scalar D, such as the interior quadratic B-spline moment
$\mat D=(\Delta x^2/4)\mat I$, but CK moments need not be scalar matrices.

\textbf{Audit boundary.} The paper-to-code document records what the code
actually does. This derivation records what follows under our stated
column-vector/gradient conventions. Resolving their difference needs an
explicit tensor-index and force-path audit; agreement between references that
share the source ordering is insufficient. This guide does not silently swap
the code's factors or claim that the discrepancy is already resolved.

Another independent distinction is paper Eq.~(17): its printed velocity
transfer includes grid-mass/particle-mass factors, whereas our shader follows
the author's CUDA average $\mat v_p=\langle\mat u\rangle$.
That source comparison is not a derivation of equality to the printed equation.
\checkpoint{The chain from Nanson through weak force to MLS is now complete.
Which multiplication order does that chain give, and which does the source
use? Keeping both answers visible is part of understanding the method.}

\section{SVD, rotation and logarithmic stretch}\label{sec:svd}
\prereq{Sections~\ref{sec:jacobian} and~\ref{sec:deformation}.}
For an invertible deformation with positive determinant, $\mat F^T\mat F$ is
symmetric positive definite. Its orthonormal eigenvectors form V and its
positive eigenvalues are $s_a^2$. Define $\mat S=\diag(s_1,s_2,s_3)$ and
$\mat U=\mat F\mat V\mat S^{-1}$. Then
\[
\mat U^T\mat U=\mat S^{-1}\mat V^T\mat F^T\mat F\mat V\mat S^{-1}=\mat I,
\qquad \boxed{\mat F=\mat U\mat S\mat V^T}.
\]
This constructs the singular value decomposition from the spectral theorem for
symmetric matrices. The values $s_a$ are principal stretches, not entries of F.
The polar rotation is $\mat R=\mat U\mat V^T$ and the right stretch is
$\mat V\mat S\mat V^T$. A rigid rotation has all $s_a=1$.

For sequential one-dimensional stretches, multiplying stretch factors should
add strain increments: $e(ab)=e(a)+e(b)$. With continuity and the small-strain
normalization $e'(1)=1$, this motivates $e(s)=\log s$. Therefore define
principal logarithmic strains $\epsilon_a=\log s_a$.
Their trace is $\sum_a\epsilon_a=\log\prod_as_a=\log J$ when F is the total
positive-determinant deformation.

Split them into mean and shape-changing parts:
\begin{equation}\label{eq:dev-strain}
t=\epsilon_1+\epsilon_2+\epsilon_3,\quad
\overline{\boldsymbol\epsilon}=\boldsymbol\epsilon-\tfrac t3\mat1,
\quad \sum_a\overline\epsilon_a=0.
\end{equation}
The vector $\mat1=(1,1,1)^T$ here is not the identity matrix.
The trace describes volume; differences among the strains describe distortion.

\textbf{Worked check.} For $\mat F=\diag(2,1,.5)$,
$\boldsymbol\epsilon=(\log2,0,-\log2)$ and $t=0$. It is volume-preserving but
not strain-free. Setting strain equal to J alone would lose its shear energy.
\checkpoint{SVD is used to expose stretches for the material model. It is not
a separate particle-grid interaction rule.}

\section{Deriving elastic stress from Hencky energy}\label{sec:energy}
\prereq{Sections~\ref{sec:stress} and~\ref{sec:svd}.}
A constitutive model specifies how deformation produces stress; MPM itself
does not determine this law. For the elastic stretch energy used in the sand
model, write an isotropic Hencky energy density
\begin{equation}\label{eq:hencky}
\Psi(\mat F)=\mu\sum_a\epsilon_a^2+\frac\lambda2
                 \left(\sum_a\epsilon_a\right)^2.
\end{equation}
Stress is the energy derivative defined by
$\delta\Psi=\mat P:\delta\mat F$.
For a singular-value energy, the derivative has the principal form
$\mat P=\mat U\diag(\partial\Psi/\partial s_a)\mat V^T$.
Since $\partial\epsilon_a/\partial s_a=1/s_a$, ordinary differentiation gives
\begin{align}
\frac{\partial\Psi}{\partial s_a}&=
    \frac{2\mu\epsilon_a+\lambda t}{s_a},\\
\mat P&=\mat U\diag\!\left(\frac{2\mu\epsilon_a+\lambda t}{s_a}\right)\mat V^T,
\label{eq:hencky-P}\\
\mat T=\mat P\mat F^T
&=\mat U\diag(2\mu\epsilon_a+\lambda t)\mat U^T.\label{eq:hencky-T}
\end{align}
For distinct stretches the principal derivative follows by differentiating
$\mat F=\mat U\mat S\mat V^T$: rotational variations do not change an isotropic
energy, leaving its diagonal stretch variation. Repeated positive stretches
are interpreted by continuity; this is not permission to differentiate a
particular numerical eigenvector sign convention.

These equations explain the \code{pHat}, U, V and final multiplication by
\code{transpose(F)} in \code{sandMaterial()}. With $s_a=1$, every logarithm is
zero and the stress is zero even for a rigid rotation.

\subsection{Why E and Poisson ratio lead to the Lam\'e parameters}
For small strains the same law linearizes to
$\boldsymbol\sigma=2\mu\boldsymbol\varepsilon
+\lambda\tr(\boldsymbol\varepsilon)\mat I$.
In uniaxial tension, take strains $(a,b,b)$ and zero lateral stresses. Then
$2\mu b+\lambda(a+2b)=0$, so
\[
\nu=-b/a=\frac{\lambda}{2(\lambda+\mu)},\qquad
E=\sigma_{11}/a=\frac{\mu(3\lambda+2\mu)}{\lambda+\mu}.
\]
Solving these two relations gives
\[
\mu=\frac{E}{2(1+\nu)},\qquad
\lambda=\frac{E\nu}{(1+\nu)(1-2\nu)}.
\]
These expressions are used in Section~\ref{sec:time} and in
\code{parameters()} in \src{src/simulation_config.cpp}.
\checkpoint{Doubling E at fixed Poisson ratio doubles both elastic moduli;
wave speed grows approximately by $\sqrt2$, reducing the acoustic step scale.}

\section{Elastic/plastic deformation and the sand yield cone}\label{sec:plastic}
\prereq{Sections~\ref{sec:svd} and~\ref{sec:energy}.}
Elastic deformation can recover; plastic deformation changes the stress-free
configuration. A common multiplicative description is
\[
\mat F^{\rm total}=\mat F^e\mat F^p.
\]
A trial update first holds plastic history fixed and evolves the elastic part:
$\mat F^{e,\rm trial}=(\mat I+\Delta t\,\mat L)\mat F^{e,n}$.
Return mapping then changes elastic stretches to satisfy a yield criterion.
This is why the code updates F in G2P and corrects it in the material pass.

The stored corrected matrix is an elastic state. Its determinant is not
necessarily the total geometric J used in Nanson's formula and density
conversion. The full plastic matrix need not be stored if the chosen stress
update only needs the elastic state and selected history variables.

Under the specific assumption that reference energy is $\Psi(\mat F^e)$ with
$\mat F^p$ held fixed during elastic differentiation,
\[
\delta\mat F^e=\delta\mat F^{\rm total}(\mat F^p)^{-1},\qquad
\mat P^{\rm total}=\mat P^e(\mat F^p)^{-T},\qquad
\mat P^{\rm total}(\mat F^{\rm total})^T=\mat P^e(\mat F^e)^T.
\]
This explains how an elastic-stretch stress calculation can supply the
Kirchhoff-like product used by the force path. Different plastic energy/volume
conventions require their own constitutive derivation.

\subsection{Pressure-dependent resistance to shear}
Use tensile stress as positive. The compression-positive pressure is
$p=-\tr\mat T/3$, and the deviatoric stress is
$\mat T_{\rm dev}=\mat T-(\tr\mat T/3)\mat I$.
A Drucker--Prager cone in the normalization used here is
\begin{equation}\label{eq:DP}
f(\mat T)=\|\mat T_{\rm dev}\|_F+\alpha\tr\mat T\le0.
\end{equation}
Equivalently the allowed shear norm is at most $3\alpha p$: compression raises
shear resistance. At zero pressure a zero-cohesion cone supports no deviatoric
stress. The implementation uses the friction-angle calibration
\[
\alpha=\sqrt{\frac23}\frac{2\sin\varphi}{3-\sin\varphi}.
\]
This calibration is a material-model convention, not a consequence of the MPM
transfer kernel. Friction-angle calibrations vary with the chosen invariant
normalization and fitting to a friction model. An internal material angle is
not a universal sliding threshold for every object or boundary.

From \eqref{eq:hencky-T}, the stress eigenvalues are $2\mu\epsilon_a+\lambda t$.
Thus $\|\mat T_{\rm dev}\|_F=2\mu\|\overline{\boldsymbol\epsilon}\|$ and
$\tr\mat T=(2\mu+3\lambda)t$. Dividing the yield function by $2\mu$ gives
\begin{equation}\label{eq:strain-yield}
\Delta\gamma=\|\overline{\boldsymbol\epsilon}\|
+\frac{3\lambda+2\mu}{2\mu}\alpha t.
\end{equation}
This derives the scalar expression seen in the return mapping when the
auxiliary volume-history correction is zero.

\subsection{Radial correction at compressive trace}
For $t<0$ and $\Delta\gamma>0$, use a non-associated, trace-preserving plastic
flow in the deviatoric direction:
\begin{equation}\label{eq:radial}
\boldsymbol\epsilon^{\rm new}
=\boldsymbol\epsilon-\frac{\Delta\gamma}
 {\|\overline{\boldsymbol\epsilon}\|}\overline{\boldsymbol\epsilon}.
\end{equation}
The trace stays t. The new deviatoric norm is
$\|\overline{\boldsymbol\epsilon}\|-\Delta\gamma
=-(3\lambda+2\mu)\alpha t/(2\mu)$, which is nonnegative under compression.
Substitution into \eqref{eq:strain-yield} gives zero: the returned state lies
on the cone. If $\Delta\gamma\le0$, the trial state is already admissible.
For tensile/noncompressive trace, the source's zero-cohesion apex treatment
sets elastic logarithmic stretches to zero. Then reconstruct
\[
s_a^{\rm new}=\exp(\epsilon_a^{\rm new}),\qquad
\mat F^{e,\rm new}=\mat U\diag(s_a^{\rm new})\mat V^T.
\]
Zero logarithmic stretch gives a rotation $\mat U\mat V^T$, not necessarily
an identity matrix. The general sand-model background is Kl\'ar et al.~\cite{sand}.
\checkpoint{Return mapping is a projection of material state, not the G2P
transfer. The transfer provides a trial deformation; plasticity decides which
part of that deformation can remain elastic.}

\section{The implemented volume-history correction}\label{sec:history}
\prereq{Section~\ref{sec:plastic}; reading a branch condition.}
The CK-MPM material overload adds a scalar history state named \code{logJp}.
The shader uses this specific algorithm, with h denoting the old history:
\begin{align*}
\epsilon_a&=\log\max(|s_a|,10^{-4}),&
s&=\sum_a\epsilon_a,& t_c&=s+h,\\
\mat e&=\boldsymbol\epsilon-\tfrac{t_c}3\mat1,&
q&=\|\mat e\|,& k&=\frac{3\lambda+2\mu}{2\mu}.
\end{align*}
If $t_c\ge0$, it sets $s_a^{\rm new}=1$ and updates $h^{\rm new}=h+s$.
If $t_c<0$, it resets $h^{\rm new}=0$ and computes
\[
\Delta\gamma=q+k\alpha t_c,\qquad
\mat H=\begin{cases}
\boldsymbol\epsilon-(\Delta\gamma/\max(q,10^{-20}))\mat e,&\Delta\gamma>0,\\
\boldsymbol\epsilon,&\Delta\gamma\le0,
\end{cases}
\qquad s_a^{\rm new}=e^{H_a}.
\]
The corrected elastic matrix and its stress are then stored.

\textbf{Important distinction.} For nonzero old h,
$\sum_a e_a=-h$, so e is not the ordinary traceless deviator from
\eqref{eq:dev-strain}. The radial-return proof above establishes the zero-history
case; it must not be presented as a proof of every volume-correction branch.
This extension is traced directly to the author's
\src{../CK-MPM/include/mpm_material.cuh} and implemented in
\src{shaders/material.glsl}. A full physical derivation of that particular
history rule requires the material-model conventions beyond the CK paper's
kernel derivation. The guide states the algorithm rather than inventing one.

Likewise, singular-value floors and denominator floors are numerical guards,
not material laws. GPU finiteness checks do not prove positive total J,
physical accuracy, constitutive admissibility in every corner case or stability
of an arbitrary manually selected timestep.
\checkpoint{When h is zero, e is the ordinary deviator and the cone derivation
applies. When h is nonzero, identify exactly which quantity changed before
reusing that proof.}

\section{Trace the dependencies through one substep}\label{sec:trace}
\prereq{All preceding branches; use the dependency map to revisit a missing link.}
The runtime dispatches
\[
\text{clear}\to\text{P2G}\to\text{grid update}\to\text{G2P}\to\text{material update}.
\]
At the beginning of a step, each particle has position, velocity, elastic F,
affine C, material stress and history. Corrected stress from the previous step
feeds the current force scatter. The initial F is identity with zero stress.

\begin{longtable}{>{\raggedright\arraybackslash}p{.20\linewidth}>{\raggedright\arraybackslash}p{.71\linewidth}}
\toprule Stage & Mathematical ancestry and implementation \\\midrule\endhead
Weights & Derivative smoothness $\to$ compact sine kernel $\to$ dual-grid moment consistency; \src{shaders/interpolation.glsl}. \\
P2G & Particle quadrature plus affine least squares supplies mass and local momentum. Nanson $\to$ Piola stress $\to$ weak force $\to$ moment-based force approximation supplies the stress term, subject to Section~\ref{sec:order-audit}; \src{shaders/p2g.comp}. \\
Grid update & Lumped mass $\to$ independent velocity normalization; explicit integration adds gravity and wall constraints; \src{shaders/grid_update.comp}. \\
G2P & Dual-grid moment identities $\to$ B and D $\to$ fitted C; chain rule $\dot F=LF$ $\to$ elastic trial update; updated velocity $\to$ advection; \src{shaders/g2p.comp}. \\
Material & SVD $\to$ logarithmic stretch $\to$ Hencky energy derivative $\to$ yield cone and return correction $\to$ stored stress/history; \src{shaders/material_update.comp}, \src{shaders/material.glsl}. \\
Checks & Numerical diagnostics summarize finite state, domain/stencil validity and maximum speed; \src{shaders/particle_health.comp}. \\\bottomrule
\end{longtable}

\subsection{A worked symbolic path to the weak nodal force}
The complete bridge that often gets compressed into one line is
\begin{align*}
\mat x=\boldsymbol\phi(\mat X)
&\Rightarrow \mat F=\nabla_X\boldsymbol\phi,\ J=\det\mat F,\\
(\mat F\mat a)\times(\mat F\mat b)
&\Rightarrow \mat n\,\mathrm da=J\mat F^{-T}\mat N\,\mathrm dA,\\
\mathrm d\mat f=\boldsymbol\sigma\mat n\,\mathrm da
&\Rightarrow \mat P=J\boldsymbol\sigma\mat F^{-T},\\
\rho_R\ddot{\boldsymbol\phi}=\Div\mat P+\rho_R\mat g
&\Rightarrow \mat f_i^{\rm int}=-\int\mat P\nabla_X w_i\,\mathrm dV,\\
\nabla_X w_i=\mat F^T\nabla_x w_i
&\Rightarrow \mat f_i^{\rm int}=-\int\mat P\mat F^T\nabla_x w_i\,\mathrm dV,\\
\int q\,\mathrm dV\approx\sum_pV_{p,0}q_p
&\Rightarrow \mat f_i^{\rm int}\approx
 -\sum_p V_{p,0}\mat P_p\mat F_p^T\nabla_x w_{i,p}.
\end{align*}
Every arrow now has a derivation above. CK changes the interpolation and dual-grid
consistency; affine/MLS constructions change the transfer and gradient
approximation. A constitutive law supplies P. These are connected but distinct
pieces of the method.

\subsection{Questions you should now be able to answer}
\begin{enumerate}
\item Why are F, J, P and $\boldsymbol\sigma$ different mathematical objects?
\item Why does surface area require $J\mat F^{-T}$ rather than J alone?
\item Which two transformations explain $\mat P\mat F^T$ in the weak force?
\item Why do the compact weights need two shifted grids for linear reproduction?
\item Why is $\mat C=\mat B\mat D^{-1}$ a fit rather than an arbitrary extra term?
\item Which factors cancel when code accumulates unaveraged cell-unit moments?
\item Which matrix ordering follows from the conventional MLS force derivation,
      and which appears in the inspected source?
\item Why does changing dt alter stability, while changing step count primarily
      changes work per rendered frame?
\item Why does sand need both an elastic stress law and a plastic correction?
\item Why is a finite returned state not itself a proof of physical correctness?
\end{enumerate}

\section{Primary references and further study}
The displayed calculations are explanatory derivations under stated assumptions.
For the missing continuum bridge, start with the Nanson/stress references; for
material-model details beyond the zero-history proof, follow the sand paper
and the exact author overload cited above. The project reference tests check
implementation behavior; they do not replace these mathematical arguments.

\begin{thebibliography}{9}
\raggedright
\bibitem{ck} M. Liu, X. Wang and M. Li.
\emph{CK-MPM: A Compact-Kernel Material Point Method}. ACM TOG 44(4), 2025.
\href{https://doi.org/10.1145/3731155}{doi:10.1145/3731155}.
Local copy: \code{../../../CK-MPM Paper.pdf}.
\bibitem{sandia} Sandia National Laboratories.
\emph{Sierra/SM Theory Manual}, Chapter 7, Stress Measures, especially
Nanson's Formula and First Piola--Kirchhoff Stress.
\url{https://www.sandia.gov/files/sierra/SM_Theory_5_18/main/stress_measures.html}.
\bibitem{nptel} NPTEL, Continuum Mechanics, Module 3, Lecture 19.
\emph{Area and Volume Transformation}.
\url{https://archive.nptel.ac.in/content/storage2/courses/112103167/module3/lec19.pdf}.
\bibitem{alberta} University of Alberta, Engineering Courses.
\emph{First and Second Piola Kirchhoff Stress Tensors}.
\url{https://engcourses-uofa.ca/books/introduction-to-solid-mechanics/stress/first-and-second-piola-kirchhoff-stress-tensors/}.
\bibitem{apic} C. Jiang et al.
\emph{The Affine Particle-In-Cell Method}. ACM TOG, 2015.
\href{https://doi.org/10.1145/2766996}{doi:10.1145/2766996}.
\bibitem{mls} Y. Hu et al.
\emph{A Moving Least Squares Material Point Method with Displacement
Discontinuity and Two-Way Rigid Body Coupling}. ACM TOG, 2018.
\href{https://doi.org/10.1145/3197517.3201293}{doi:10.1145/3197517.3201293}.
\bibitem{sand} G. Kl\'ar et al.
\emph{Drucker--Prager Elastoplasticity for Sand Animation}. ACM TOG, 2016.
\href{https://doi.org/10.1145/2897824.2925906}{doi:10.1145/2897824.2925906}.
\end{thebibliography}
\end{document}
